\section*{第 \ref{ch:b3:genomics} 章——基因组学与测序}

\begin{solution}{exo:b3:genomics:1}
\omterm{def:b3:genomics:sanger}{双脱氧核苷酸}没有 3$'$ 羟基，因此无法
与下一个核苷酸形成磷酸二酯键，链随之终止。若只有双脱氧形式，每条链
都会在第一个位置终止；若只有普通形式，
则没有链会终止。混合使用时，每个位置上的终止都是随机事件，
于是产物形成完整的梯状图谱，
模板的每个碱基各对应一种长度。
\end{solution}

\begin{solution}{exo:b3:genomics:2}
$4\times 10^{8}\times \qty{300}{bp} = 1.2\times 10^{11}$ 个碱基。人类：
$1.2\times 10^{11}/3.2\times 10^{9} \approx 38\times$。细菌：
$1.2\times 10^{11}/5\times 10^{6} = \num{24000}\times$。
\end{solution}

\begin{solution}{exo:b3:genomics:3}
重叠群是由重叠\omterm{def:b3:genomics:ngs}{读段}拼接而成的连续序列；\omterm{def:b3:genomics:assembly}{支架}是按顺序、定向排列的
一组重叠群，相邻重叠群之间隔以估计大小的空位；\omterm{def:b3:genomics:assembly}{N50} 是将重叠群
按长度从大到小排序后，累计碱基数达到\omterm{def:b3:genomics:assembly}{组装}碱基总数一半时的重叠群长度。
这里十个 \qty{8}{Mb} 的重叠群合计含 \qty{80}{Mb}，超过 \qty{50}{Mb} 的一半
（累计到第七个时达到 \qty{56}{Mb}），
故 \omterm{def:b3:genomics:assembly}{N50} $= \qty{8}{Mb}$。
\end{solution}

\begin{solution}{exo:b3:genomics:4}
\omterm{prop:b3:genomics:sizes}{基因组大小}并不随生物体的复杂程度而变化：洋葱的 DNA
含量是人类的五倍，肺鱼则是其四十倍；酵母与线虫的 DNA
含量相差十倍，而基因数目相近。多出的部分是非编码序列，
包括转座元件及其残迹、内含子、
卫星重复序列。
\end{solution}

\begin{solution}{exo:b3:genomics:5}
$e^{-c} = 10^{-6}$ 给出 $c = 6\ln 10 = 13.8$。取 $c = 8$，则 $N =
cG/L = 8\times 5\times 10^{6}/150 \approx \num{267000}$ 条\omterm{def:b3:genomics:ngs}{读段}，$\theta
= 30/150 = 0.2$，重叠群数 $= N e^{-6.4} = \num{267000}\times 0.00166
\approx 440$。
\end{solution}

\begin{solution}{exo:b3:genomics:6}
$P(X < 8) = P(X \le 7) \approx P\bigl(Z < (7.5 - 15)/2.74\bigr) =
P(Z < -2.74) \approx 0.003$。在 $30\times$ 覆盖度下，杂合子的两个等位基因
几乎总能被多次观测到，覆盖度分布不均
（富含 GC 的区域\omterm{def:b3:genomics:ngs}{读段}较少，部分位点的覆盖度只有平均值的一半），
而\omterm{def:b3:genomics:ngs}{读段}数仍足以将真实等位基因与 $10^{-3}$
的误差区分开来。
\end{solution}

\begin{solution}{exo:b3:genomics:7}
重复序列内部的每条 \qty{150}{bp} \omterm{def:b3:genomics:ngs}{读段}，无论来自哪个拷贝，其序列都完全相同，
因此\omterm{def:b3:genomics:assembly}{组装}程序构建了一个 \qty{6}{kb} 的节点，该节点有
$500$ 条路径进入、$500$ 条路径离开；它无法判断哪个入口
与哪个出口配对，\omterm{def:b3:genomics:assembly}{组装}因此在 $500$ 处断开，重复序列只出现一次，
但其覆盖度为平均值的 $500$ 倍。只有长于重复序列加两侧独特侧翼的\omterm{def:b3:genomics:ngs}{读段}
（即 \qty{8}{kb} 或更长）才能跨越每个拷贝
并将其分辨。
\end{solution}

\begin{solution}{exo:b3:genomics:8}
$4.8\times 10^{7}/1000 = \num{48000}$ 个编码区差异；其中三分之一，
约 \num{16000} 个会改变蛋白质。
\end{solution}

\begin{solution}{exo:b3:genomics:9}
期望的假阳性数为 $10^{6}\times 5\times 10^{-8} = 0.05$。对于
患病率为 \qty{2}{\%} 的疾病，相对风险为 $1.15$ 时，病例数仅变化几例，
即每千人中仅几例；千人中约含二十例病例，远不足以
察觉这一效应（检验效能随病例数及效应量的平方增长）。
对个体而言，$0.3$ 个百分点毫无意义。
但该等位基因所标记的基因，其活性的微小改变
会改变患病风险——该基因及其通路可能是一个
药物靶点，完全抑制它会产生较大效应。
\end{solution}

\begin{solution}{exo:b3:genomics:10}
群体遗传角度：细菌种群规模巨大，选择
会清除即使代价很小的插入；哺乳动物的有效种群较小，
漂变使轻度有害的非编码 DNA 得以积累——
这一点得到各谱系中\omterm{prop:b3:genomics:sizes}{基因组大小}与种群大小
呈负相关的支持，但也受到例外的削弱。
结构角度：真核生物基因组被转座子入侵，而细菌
借助缺失偏向，且自私元件没有减数分裂的庇护，将其清除——
该解释得到\omterm{prop:b3:genomics:sizes}{基因组大小}
与转座子含量相关性的支持。调控角度：复杂发育需要更多
调控 DNA——这一点得到\omterm{def:b3:genomics:comparative}{保守非编码元件}的支持，
即发育基因周围的保守元件；但事实上只有
\qty{8}{\%} 的人类基因组显示出任何选择信号，因此大多数
非编码 DNA 并非调控性的。
\end{solution}

\begin{solution}{exo:b3:genomics:11}
两类起点以总密度 $\rho = (N_{1} +
N_{2})/G$。长度为 $L_{i}$ 的\omterm{def:b3:genomics:ngs}{读段}位于其重叠群最右端，当且仅当
其后 $L_{i} - T$ 个碱基内没有任何一类\omterm{def:b3:genomics:ngs}{读段}起始：
概率为 $e^{-\rho(L_{i} - T)}$。由此得重叠群数 $= N_{1}
e^{-\rho(L_{1} - T)} + N_{2} e^{-\rho(L_{2} - T)}$。长\omterm{def:b3:genomics:ngs}{读段}
位于最右端的概率按指数规律小于短\omterm{def:b3:genomics:ngs}{读段}，
且每条长\omterm{def:b3:genomics:ngs}{读段}还消除了其整个长度上出现空位的机会；
与短\omterm{def:b3:genomics:ngs}{读段}碱基数相同时，会增加许多起始点，但每个受保护的窗口较短，
故长\omterm{def:b3:genomics:ngs}{读段}更占优势。
\end{solution}

\begin{solution}{exo:b3:genomics:12}
两次\omterm{def:b3:genomics:comparative}{全基因组复制}预测，这一四重模式在全基因组范围内
成立：四联的旁系同源染色体片段（旁系同源区段）
携带相同的基因家族，且顺序相同，复制时间
在所有家族中均定于同一时刻；仅有一个簇，
见于文昌鱼与\emph{Ciona}，二者在这些事件之前已分化；以及
七鳃鳗则在这些事件前后分化，属于中间或
独立衍生状态。Hox 簇的独立复制
仅预测 Hox 有旁系同源区段，邻近基因的
历史各不相同，各家族的复制时间也各异。
全基因组范围的旁系同源区段与一致的复制定年，
与观察结果相符，支持\omterm{def:b3:genomics:comparative}{全基因组复制}。
\end{solution}

\begin{solution}{pb:b3:genomics:1}
\textbf{1.} $0.002\times 6\times 10^{8} = 1.2\times 10^{6}$ 条\omterm{def:b3:genomics:ngs}{读段}；
$c = 1.2\times 10^{6}\times 150/4.6\times 10^{6} \approx 39$.
\textbf{2.} $e^{-39} \approx 10^{-17}$：实际上没有碱基未被覆盖。
\textbf{3.} $\theta = 0.2$；重叠群数 $= 1.2\times 10^{6} e^{-31} \approx
0$：理论上只有一个重叠群，空位仅出现在重复序列处。
\textbf{4.} $c = \num{30000}\times 150/4.6\times 10^{6} = 0.98$；
未覆盖部分 $e^{-0.98} = 0.38$，约 \qty{1.7}{Mb}；重叠群数 $= \num{30000}
\times e^{-0.78} \approx \num{13700}$.
\textbf{5.} 七个操纵子给出相同的\omterm{def:b3:genomics:ngs}{读段}，并塌缩为一个
\qty{5}{kb} 重叠群，由七个各不相同的左侧侧翼进入，
由七个右侧侧翼离开（其配对关系未知）：$14$ 个重叠群末端，七个
空位。
\textbf{6.} 约 \num{4600} 个基因；编码序列 $0.88\times 4.6\times 10^{6}
= 4.0\times 10^{6}$ bp，平均基因长 \qty{880}{bp}。
\textbf{7.} $0.998\times 6\times 10^{8}\times 150/3.2\times 10^{9}
\approx 28\times$.
\textbf{8.} 每个等位基因约 $14$ 条\omterm{def:b3:genomics:ngs}{读段}。
\textbf{9.} $28\times 10^{-3}/3 \approx 0.009$ 条\omterm{def:b3:genomics:ngs}{读段}显示某一给定的错误碱基——
三条或更多的概率约为 $10^{-7}$——且
\qty{20}{\%} 的 $28$ 即 $5.6$ 条\omterm{def:b3:genomics:ngs}{读段}；此错误两项检验均不通过，
而真实的杂合等位基因预期出现在 $14$ 条\omterm{def:b3:genomics:ngs}{读段}中。
\textbf{10.} $3.2\times 10^{9}/1500 \approx 2.1\times 10^{6}$
杂合位点数；编码序列中为 $4.8\times 10^{7}/1500 \approx
\num{32000}$.
\textbf{11.} 来自 Alu 的 \qty{150}{bp} \omterm{def:b3:genomics:ngs}{读段}与许多 $1.1$
百万个拷贝几乎同样匹配，因此它被定位到错误的拷贝上，
或得到较低的比对置信度。此时拷贝之间的差异
看起来像杂合变异，而真实变异则被隐藏其中
：Alu 内的变异检出通常被过滤掉，而
这些元件是短读长测序的盲区。
\textbf{12.} $1.2\times 10^{-8}\times 6.4\times 10^{9} \approx 77$ 个新
突变。子代中该变异应在约 $14$ 条\omterm{def:b3:genomics:ngs}{读段}中出现；但
若某亲本位点仅被五条\omterm{def:b3:genomics:ngs}{读段}覆盖，则杂合等位基因
被漏检的概率为 $2^{-5} = 3\,\%$，而遗传来的变异则
被误判为新突变——因此亲本的测序深度必须与
子代相当。
\textbf{13.} $4.8\times 10^{7}\times 100/150 = 3.2\times 10^{7}$ 条\omterm{def:b3:genomics:ngs}{读段}，
比全基因组少二十倍：测序成本
降低了二十倍，超过捕获步骤的成本。
\textbf{14.} $1.1\times 10^{6}\times 300 = 3.3\times 10^{8}$ bp，约占
基因组的 \qty{10}{\%}；\qty{10}{\%} 的\omterm{def:b3:genomics:ngs}{读段}，约 $6\times
10^{7}$.
\textbf{15.} $\num{60000}\times 1500 = 9\times 10^{7}$ bp；$9\times
10^{7}/1.6\times 10^{10} = 0.56\,\%$.
\textbf{16.} $0.012\times 2.9\times 10^{9} = 3.5\times 10^{7}$
个替换，每个谱系 $1.7\times 10^{7}$；速率 $1.7\times
10^{7}/(2.9\times 10^{9}\times 6.5\times 10^{6}) = 9\times 10^{-10}$
每碱基每年；以 25 年为一代时为 $2.3\times 10^{-8}$——
约为系谱速率 $1.2\times 10^{-8}$ 的两倍，这是一个已知的
差异，提示过去的世代时间更长，或
分歧发生得更早。
\textbf{17.} 若四份各含 \num{16000} 个基因的拷贝，应为 \num{64000}；现有 \num{20000}
个保留，因此在 \num{48000} 个额外拷贝中，仅 \num{4000} 个留存：
\qty{92}{\%} 的重复拷贝丢失了。
\textbf{18.} 线虫与人类的基因数目相同；复杂性
在于基因的利用方式——可变剪接可使一个基因
产生数千种蛋白质（\emph{Dscam} 可达 \num{38016} 种），调控
将转录因子与增强子以细胞类型特异的方式
组合；非编码 RNA 增加了调控层次，蛋白质则在网络中相互作用。
\textbf{19.} 调控元件（增强子、启动子、绝缘子），
非编码 RNA 基因、剪接与定位信号、复制起点和
结构序列。检验一个增强子候选序列，可将其置于
报告基因上游，导入转基因胚胎中，观察报告基因在何处
表达，然后借助 CRISPR 删除基因组中的该元件，
并测定邻近基因的表达量。
\textbf{20.} $0.05/10^{6} = 5\times 10^{-8}$；在 $p < 10^{-5}$ 时，$10^{6}
\times 10^{-5} = 10$ 个预期假阳性。
\textbf{21.} 非携带者的比数为 $0.009/0.991 = 0.00908$；单拷贝比数
为 $0.00954$，风险 $0.945\,\%$；双拷贝比数 $0.01001$，风险
$0.99\,\%$.
\textbf{22.} 该得分对 $200$ 个等位基因求和，即每人携带的风险
拷贝数，以各等位基因的对数比值比加权。其
计数的均值为 $200\times 2\times 0.3 = 120$，标准差
为 $\sqrt{200\times 2\times 0.3\times 0.7} \approx 9$；排名前 \qty{1}{\%} 的人
平均多携带约 $21$ 个等位基因，且 $1.05^{21} \approx
2.8$：比数几乎为普通人的三倍，这些等位基因
各自仅使风险变化 \qty{5}{\%}。
\textbf{23.} 位于同一数十千碱基区段内的等位基因会一起遗传
（连锁不平衡），因此其中任一位点都显示相同的
关联；被分型的 SNP 只是芯片上的那一个。要找到
因果变异：在病例与对照中对该区域测序，缩小
统计上的候选集，再逐一检验：报告基因实验
检测各等位基因的增强子活性，在细胞中编辑各等位基因并
测定邻近基因的表达，并
与表达数量性状信号进行共定位。
\textbf{24.} $1 - e^{-0.5} = 39\,\%$ 的基因组被读取。这提供了
某一已知过去时间点种群的直接样本——等位基因
频率（早于后来的迁移与混合），以及
尼安德特人片段的长度（较长，因为重组代数少，
未曾打断这些片段），并由此得到混合的时间——这是现代
基因组只能借助模型推断的。
\textbf{25.} 试点中约有 \num{13700} 个重叠群；每个单倍体基因组的覆盖深度为 $28\times$，
每个等位基因约 $14$ 条\omterm{def:b3:genomics:ngs}{读段}；子代中约有 $77$ 个
新突变。
\end{solution}
